\documentclass[10pt]{article}

\usepackage[utf8]{inputenc}

\usepackage[T1]{fontenc}

\usepackage{amsmath}

\usepackage{amsfonts}

\usepackage{amssymb}

\usepackage[version=4]{mhchem}

\usepackage{stmaryrd}

\usepackage{graphicx}

\usepackage[export]{adjustbox}

\graphicspath{ {./images/} }

\usepackage{bbold}



\DeclareUnicodeCharacter{0131}{$\imath$}



\begin{document}

ЕсліІ $X$ и $\boldsymbol{Y}$-случайвые величины, $f_{Y}(y)$-плотность распределения $Y, f_{X}, Y(x, y)$-совместная плотность распределения $X$ и $Y$, то



\[

f_{X}(x \mid Y=y)=\frac{1}{f_{Y}^{(y)}} f_{X, Y}(x, y)

\]



определяет У. п. распределения случайной величины $X$ гри фиксированном значении $Y$.



\includegraphics[max width=\textwidth, center]{2023_07_09_2c64cdb37c269824a52fg-1}

УСЛОВНАЯ СХОДИМӦСТЬ р я а а - свойство ряда, заключающееся в том, что существует сходящийся ряд, полученный из данного нек-рой перестановкой его членов. Числової ряд



\[

\sum_{n=1}^{\infty} u_{n}

\]



б е зусло вно сх о ди т с я, если он сходится, и сходится любоӥ ряд, полученный перестановкой его членов, причем сумма любого такого ряда одна и та же, иначе говоря, сумма безусловно сходящегося ряда нс зависит от порядка его членов. Если ряд (\textit{) сходится, но не безусловно, то он наз. у с ло вно с ходящ и мся. Для того чтобы ряд (}) условно сходился, необходимо и достаточно, чтобы он сходился, но не абсолютно, т. е. чтобы $\sum_{n=1}^{\infty}\left|u_{n}\right|=+\infty$.



Если члены ряда (\textit{) являются действительными числами, через $u_{1}^{+}, u_{2}^{+}, \ldots, u_{n}^{+}, \ldots$ обозначены его неотрицательные члены, а через - $u_{\mathbf{1}}^{-}, \quad-u_{2}^{-}, \ldots$ $\ldots,-u_{n}, \ldots$-отрицательные, то ряд (}) будет условно сходиться тогда ІІ только тогда, когда оба ряда $\sum_{n=1}^{\infty} u_{n}^{+}$ı $\sum_{n=1}^{\infty} u_{n}^{-}$расходятся (при этом порядок слагаемых в этих рядах безразличен).



Пусть ряд (\textit{) с действительными членами сходится условно и $-\infty \leqslant \alpha<\beta \leqslant+\infty$, тогда существует такой ряд $\sum_{n=1}^{\infty} u_{n}^{*}$, полученный перестановкой членов ряда (}), что если обовначить через $\left\{s_{n}^{*}\right\}$ последовательвость его частичных сумм, то



\[

\varliminf_{n \rightarrow \infty} s_{n}^{*}=\alpha, \varlimsup_{n \rightarrow \infty} s_{n}^{*}=\beta

\]



(это есть обобщение т е о р е м ы Р и м а в а).



П роизведение условно сходящихся рядов зависит от порядка, в к-ром суммируются результаты почленного умножения членов данных рядов.



Понятие условной и безусловной сходимости ряда обобщается на ряды, члены к-рых являются элементами нек-рого нормиронанного векторного п ространства $X$. Если $X$-конетномерное пространство, то аналогично случаю числовых рядов сходящийся ряд $\sum_{n=1}^{\infty} u_{n}, u_{n} \in X, n=1,2, \ldots$, условно сходится тогда и только тогда, когда ряд $\sum_{n=1}^{\infty}\left\|u_{n}\right\|_{X}$ расходится. Если же пространство $X$ бесконечномерное, то в нем существуют безусловно сходящиеся ряды $\sum_{n=1}^{\infty} u_{n}$, не являющиеся абсолютно сходящимися, т. с. такие, что для них $\sum_{n=1}^{\infty}\left\|u_{n}\right\|_{\boldsymbol{X}}=+\infty$.



Л. Д. Кудрявчев.



УСЛОВНАЯ УСТОЙЧВОСТЬ Точки относительно семеӥства отображенийі



\[

\left\{f_{t}\right\}_{t \in G+}: E \longrightarrow E

\]



\begin{itemize}

  \item равностепенная непрерывность в этой точке семейства $\left\{\left.f_{t}\right|_{V}\right\}_{t \in G+}$ сужений отображений $f_{t}$ на нек-рое вложенное в $E$ многообразие $V$; здесь $G+-$ множество неотрицательных чисел: действительных $(G=\mathbb{R})$ или целых $(G=\mathbb{Z})$.

\end{itemize}



У. У. точкі относительно отображения оп ределяется как ее У. У. относптельно семейства неотрицатөльных степеней әтого отображения. У. у. точки относительно динамич. системы $f^{t}$ есть У. у. этой точки относительно семейства отображений $\left\{f^{t}\right\}_{t \in G+\cdot}$ У. у. заданного на $t_{0}+\mathbb{Z}+$ решения $x_{0}(\cdot)$ уравнения



\[

x(t+1)=g_{t} x(t)

\]



есть У. У. точки $x_{0}\left(t_{0}\right)$ относительно семейства отображений



\[

\left\{f_{t}=g_{\operatorname{def}}^{=} g_{t_{0}+t} \ldots g_{t_{0}+1} g_{t_{0}}\right\}_{t \in \mathbb{Z}^{+}} .

\]



У. у. заданного на $t_{0}+\mathbb{R}+$ реппения $x_{0}(\cdot)$ дифференциального уравнения



\[

\dot{x}=f(x, t)

\]



есть У. У. точки $x_{0}\left(t_{0}\right)$ относительно семейства отображений $\left\{X\left(t_{0}+t, t_{0}\right)\right\}_{t \in \mathbb{R}}$, где $X(\theta, \tau)-$ Koши операmор этого уравнения. У. У. заданного на $t_{0}+\mathbb{R}+$ решения $y(\cdot)$ диффференциального уравнения



\[

y^{(m)}=g\left(y, \dot{y}, \ldots, y^{(m-1)}, t\right)

\]



порядка $m$ есть У. У. заданного на $t_{0}+\mathbb{R}+$ решения $x(\cdot)=\left(y(\cdot), \dot{y}(\cdot), \ldots, y^{(m-1)}(\cdot)\right)$ соответствующего дифференциального уравнения 1-го порядка вида (2), где\[

x=\left(x_{1}, \ldots, x_{m}\right) \text {, }

\][



f(x, t)=\textbackslash left(x\_\{2\}, ···, x\_\{m\}, g\textbackslash left(x\_\{1\}, ···, x\_\{m\}, t\textbackslash right)\textbackslash right) .

]



Приведенные ниже определения $1-5$ являются нек-рыми конкретизациями этих и родственных им определений.



\begin{enumerate}

  \item Пусть дано дифференциальное уравнение (2), где $x \in E, E$ - нормированное $n$-мерное п ространство. Решевие $x_{0}(\cdot): l_{0}+\mathbb{R}+\longrightarrow E$ этого уравнения наз. у с л о в в о у с т о й чи вым с инд ексом $\boldsymbol{k} \in\{0, \ldots, n\}$, если найдется вложкеный в $E k$-мерный диск $D^{k}$ (рассматриваемый как многообразие нек-рого класса $\left.C^{m}\right)$, содержащий точку $x_{0}\left(t_{0}\right)$ и обладающий свойством: для всякого $\varepsilon>0$ найдется $\delta>0$ такое, тто для всякого $\quad x \in D^{k}, \quad$ удовлетворяющего неравенству $\left|x-x_{0}\left(t_{0}\right)\right|<\delta$, решение $x(\cdot)$ того же уравнения, удовлетворяющее начальному условию $x\left(t_{0}\right)=x$, единственно, определено на $t_{0}+\mathbb{R}+$ и при всяком $t \in t_{0}+\mathbb{R}+$ удовлетворяет неравенству $\left|x(t)-x_{0}(t)\right|<\varepsilon$. Если диск $D^{k}$, Іимюццй указанные свойства, может быть выбран так, что

\end{enumerate}



\[

\lim _{t \rightarrow+\infty}\left|x(t)-x_{0}(t)\right|=0

\]



(соответственно,



\[

\varlimsup_{t \rightarrow+\infty} \frac{1}{t} \ln \left|x(t)-x_{0}(t)\right|<0

\]



здесь и далее считается, тто $\ln 0=-\infty$ ) для всякого решения того же уравнения, начинающегося в этом диске (т. е. такого, что $\left.x\left(t_{0}\right) \in D^{k}\right)$, то решение $x_{0}(t)$ наз. а с и м п тот и ч е с к и (соответственно, эк с п оненйа Јьно) Ус ловно устойчивым (с индексом $k$ ).



Решение уравнения (2) $\left(x \in \mathbb{R}^{n}\right.$ или $\left.x \in \mathbb{C}^{n}\right)$ наз. условно (асимптотически, әкспоненци а льно условно) ус тойтивым с индексом $k$, если оно становится таковым в результате наделения пространства $\mathbb{R} n$ (или $\mathbb{C}^{n}$ ) нек-рой нормой. От выбора нормы это свойство решения не зависит.



\begin{enumerate}

  \setcounter{enumi}{1}

  \item Пусть дано $n$-мерное риманово многообразие $V^{n}$ (расстояние в к-ром обоозначается через $d(\cdot, \cdot)$ ). Тотка $x_{0} \in V^{n}$ наз. ус ло вно у с тойччи во й (с индексом

\end{enumerate}



\end{document}